    \documentclass[a4paper,6pt]{extarticle}
    \usepackage[utf8]{inputenc}
    \usepackage[landscape,top=0.5cm,bottom=0.5cm,left=0.5cm,right=0.5cm]{geometry}
    \usepackage{amsmath}
    \usepackage{amssymb}
    \usepackage{multicol}
    \usepackage{enumitem}
    \usepackage{xcolor}
    \usepackage{titlesec}

    \setlength{\parindent}{0pt}
    \setlength{\parskip}{0.8pt}
    \setlength{\columnsep}{0.18cm}
    \setlist{nosep,leftmargin=*,itemsep=0pt}
    \titlespacing*{\section}{0pt}{2pt}{1pt}
    \titlespacing*{\subsection}{0pt}{1.5pt}{0.5pt}
    \titleformat*{\section}{\small\bfseries\color{blue}}
    \titleformat*{\subsection}{\footnotesize\bfseries\color{red}}

    \pagestyle{empty}

    \begin{document}

    \begin{multicols}{3}
    \footnotesize

    \section*{RF - RADIOFRECUENCIA}

\subsection*{Antenas}
\textbf{Ganancia:}
$$\boxed{G_{dBi} = 10\log_{10}(\eta) + 20\log_{10}(\pi d) - 20\log_{10}(\lambda)}$$
• $\eta$ = 0.5-0.7 • $d$ [m] • $\lambda$ [m] • $A_e = G\lambda^2/(4\pi)$

\textbf{Ancho haz:} $\boxed{\theta_{3dB} \approx 70\lambda/d}$ [°]

\subsection*{Pérdidas Espacio Libre (FSL)}
$$\boxed{A_{FSL,dB} = 92.45 + 20\log_{10}(d_{km}) + 20\log_{10}(f_{GHz})}$$
• MHz: usar 32.45 (diferencia 60 dB) • Origen: expansión esférica $1/r^2$

    \subsection*{Pérdidas Adicionales}
    $\boxed{A_{total} = A_{FSL} + A_{atm} + A_{rain} + A_{pol} + A_{point}}$
    $A_{atm} = \alpha_{atm}d$ • $A_{rain} = \alpha_R L_{eff}$ • $L_{eff} = h_R/\sin\varepsilon$
    $A_{pol} = -20\log|\cos\Delta\theta|$ • $A_{point} \approx 12(\theta_{err}/\theta_{3dB})^2$

    \subsection*{Link Budget}
    $\boxed{P_R = P_{TX} + G_{TX} + G_{RX} - A_{total}}$
    
    $C/N = P_R - N$ • $M = P_R - P_{sens}$
    
    \textbf{Multi-hop:} $\boxed{1/(C/N)_{tot} = \sum 1/(C/N)_i}$ [LINEAL]

\subsection*{Ruido}
$$\boxed{N_{dBm} = -174 + 10\log_{10}(B_{Hz}) + 10\log_{10}(T_{sys}/290)}$$
• $k = 1.38\times10^{-23}$ J/K • $T_{sys} = T_A + T_{eq}$ • $T_{eq} = T_0(F-1)$

$$\boxed{F_{total} = F_1 + \frac{F_2-1}{G_1} + \frac{F_3-1}{G_1G_2} + \cdots}$$ [LINEALES]

• Primera etapa (LNA) domina • Usar $F$ y $G$ lineales

    \subsection*{PA y Modulación}
    $P_{nom} = P_{1dB} - BO$ • $P_{IM3,out} = 3P_{in} - 2P_{IP3}$
    
    $\boxed{\eta = \log_2(M)}$ • $R_b = B\eta$ • $E_b/N_0 = (C/N)(B/R_b)$
    
    Shannon: $\boxed{C = B\log_2(1 + C/N)}$

    \subsection*{Enlaces Satelitales}
    $\boxed{\Delta f = f_0(v/c)\cos\theta}$ (Doppler)
    
    $\boxed{G/T = G_{dBi} - 10\log(T_{sys})}$ [dB/K]
    
    $\cos\varepsilon = (R_T + h)/d_{sat} - (R_T/d_{sat})\cos(DL)\cos(l)$

    \subsection*{Receptores}
    \textbf{Superheterodino:} $f_{IF} = |f_{RF} - f_{LO}|$ • $f_{image} = f_{LO} + f_{IF}$
    
    \textbf{Homodino:} $f_{LO} = f_{RF}$ • $f_{IF} = 0$

    \subsection*{Atenuación Atmosférica}
    \textbf{Gases:} $\alpha_{atm} \approx 0.01$ dB/km (< 10 GHz) • $\alpha_{O_2}$ @ 60 GHz, $\alpha_{H_2O}$ @ 22 GHz
    
    \textbf{Lluvia:} $\boxed{\alpha_R = a R^b}$ • $R$ [mm/h] • ITU-R P.838

    \subsection*{Patrones Antena}
    $D = 4\pi/\Omega_A$ • $\Omega_A = \theta_{3dB,E} \theta_{3dB,H}$ [sr]
    
    \textbf{SLL:} -13 dB (uniforme), -20 dB (Taylor)
    
    \textbf{Polarización:} Lineal (V/H), Circular (RHCP/LHCP) • AR: 0 dB (circ.), $\infty$ (lin.)

    \section*{ÓPTICA - FIBRA ÓPTICA}

    \subsection*{Line Rate}
    $\boxed{R_{line} = N_{pol} R_s \log_2(M) (1+OH_{FEC})}$
    
    $R_{payload} = R_{line}/(1+OH)$ • $B_{ch} = R_s(1+\beta)$ • $\eta = R_{payload}/B_{ch}$

\subsection*{OSNR (Optical SNR)}
$$\boxed{\text{OSNR}_{dB} = P_{ch,dBm} - 10\log_{10}(P_{ASE,tot,mW})}$$
$$\boxed{\text{OSNR}_{dB} = P_{ch,dBm} - F_{dB} - 58 - 10\log_{10}(N_{amp})}$$
• Medido en $B_{opt} = 12.5$ GHz • $F_{dB}$ = 4-6 dB • 58 dB = constante teórica

\textbf{ASE:} $P_{ASE} = 2n_{sp}h\nu GB_{opt}$ • $n_{sp}$ = 1-2 • $F = 2n_{sp}$

    \subsection*{Ingenieria Enlace}
    \textbf{Número amplificadores:}
    $$\boxed{N_{amp} = \lceil L/D_{amp}\rceil - 1}$$
    $$\boxed{D_{amp} = G_{max}/\alpha}$$
    • $L$ = longitud total enlace [km] • $\alpha$ = atenuación fibra [dB/km]
    • $D_{amp}$ = espaciado entre amplificadores [km]
    
    \textbf{Número canales WDM:}
    $$\boxed{N_{ch} = \lfloor B_{amp}/\Delta f\rfloor}$$
    $$\boxed{P_{ch} = P_{out} - 10\log_{10}(N_{ch})}$$
    • $B_{amp}$ = ancho banda EDFA (típ. 35 nm) • $\Delta f$ = espaciado ITU-T
    
    \textbf{Dispersión cromática:}
    $$\boxed{D_{tot} = D \times L}$$
    • $D \approx 17$ ps/(nm·km) @ 1550nm • $D_{tot}$ < límite modulación
    
    \textbf{Enlace no amplificado:}
    $$\boxed{L_{max} = \frac{P_{TX} - P_{RX,sens} - L_{conn} - L_{MUX}}{\alpha}}$$

    \subsection*{Compensación Dispersión}
    $\boxed{L_{DCF} = -D_{SMF} L_{SMF}/D_{DCF}}$ • $\Delta P \approx (D \cdot L \cdot \Delta\lambda)^2 / 2$

    \subsection*{WDM}
    $\boxed{N_{ch} = \lfloor B_{amp}/(\Delta f + f_{guard})\rfloor}$ • $P_{total} = P_{ch} + 10\log(N_{ch})$
    
    ITU-T: 12.5, 25, 50, 100 GHz

    \subsection*{Efectos No Lineales}
    \textbf{SPM:} Self-Phase Modulation • \textbf{XPM:} Cross-Phase Modulation
    \textbf{FWM:} $f_{ijk} = f_i + f_j - f_k$ • \textbf{SRS:} Raman • \textbf{SBS:} $P_{th} \approx 1$ mW

    \subsection*{BER y Codificación}
    \textbf{BER:} < $10^{-12}$ (sin FEC), < $10^{-3}$ (con FEC)
    
    $\boxed{Q = 20\log(\sqrt{2}\text{erfc}^{-1}(2\cdot\text{BER}))}$ [dB]
    
    BER $10^{-3} \rightarrow$ 9.8 dB • $10^{-9} \rightarrow$ 15.6 dB
    
    \textbf{FEC:} 3-6 dB • RS: OH=6.7\% • LDPC: OH=20-25\%

    \subsection*{Modelos Propagación}
    \textbf{Fresnel:} $r_1 = \sqrt{\lambda d_1 d_2/(d_1+d_2)}$ • Despeje 60\%: < 0.6$r_1$
    
    \textbf{Reflexión:} $\Gamma = (n_2-n_1)/(n_2+n_1)$

    \section*{CONVERSIONES}

    \subsection*{dB ↔ Lineal}
    $\boxed{X_{dB} = 10\log(X)}$ • $X = 10^{X_{dB}/10}$ • $P_{dBW} = P_{dBm} - 30$
    
    3dB=$\times$2 • 6dB=$\times$4 • 10dB=$\times$10 • 20dB=$\times$100

    \subsection*{Frecuencia ↔ Longitud Onda}
    $\boxed{\lambda = c/f}$ • $c = 3\times10^8$ m/s
    
    \textbf{Banda C:} 1530-1565 nm • 100GHz = 0.8nm @ 1550nm

    \subsection*{Potencia e Impedancia}
    $P = V^2/R = I^2 R$ • $R_0 = 50\Omega$ • $P_{dBm} = 10\log(P_{mW})$
    
    \textbf{VSWR:} $s = (1+|\Gamma|)/(1-|\Gamma|)$ • $|\Gamma| = (s-1)/(s+1)$
    
    $A_{mismatch} = -10\log(1-|\Gamma|^2)$ • $RL = -20\log|\Gamma|$

    \subsection*{Constantes}
    \begin{tabular}{|l|l|l|}
    \hline
    \textbf{Símbolo} & \textbf{Valor} & \textbf{Descripción} \\
    \hline
    $c$ & $3\times10^8$ m/s & Velocidad luz \\
    $k$ & $1.38\times10^{-23}$ J/K & Boltzmann \\
    & $-228.6$ dBW/K/Hz & (en dB) \\
    $h$ & $6.626\times10^{-34}$ J·s & Planck \\
    $T_0$ & 290 K & Temp. referencia \\
    $\nu_{1550}$ & 194 THz & Frec. @ 1550nm \\
    & $-174$ dBm & $kT_0$ @ 1Hz \\
    \hline
    \end{tabular}

    \subsection*{Bandas Frecuencia}
    \textbf{RF:} L(1-2) S(2-4) C(4-8) X(8-12) Ku(12-18) K(18-27) Ka(27-40) V(40-75) [GHz]
    \textbf{Óptica:} O(1260-1360) E(1360-1460) S(1460-1530) C(1530-1565) L(1565-1625) [nm]
    \textbf{Satélite:} C(6↑/4↓) Ku(14↑/11-12↓) Ka(30↑/20↓) [GHz]

    \subsection*{Márgenes Sistema}
    $\boxed{M = P_R - P_{sens}}$ • Mín: 3 dB • Rec: 6-10 dB

    \section*{TABLAS REFERENCIA}

    \subsection*{Modulaciones RF}
    \begin{tabular}{|l|c|c|}
    \hline
    \textbf{Mod} & \textbf{C/N [dB]} & \textbf{$\eta$ [b/s/Hz]} \\
    \hline
    BPSK & 6-9 & 1 \\
    QPSK & 10-13 & 2 \\
    16QAM & 16-20 & 4 \\
    64QAM & 24-28 & 6 \\
    256QAM & 32-36 & 8 \\
    \hline
    \end{tabular}

    \subsection*{Modulaciones Ópticas}
    \begin{tabular}{|l|c|}
    \hline
    \textbf{Modulación} & \textbf{OSNR [dB]} \\
    \hline
    DP-QPSK & 10-13 \\
    DP-8QAM & 15-17 \\
    DP-16QAM & 19-22 \\
    DP-64QAM & 27-30 \\
    \hline
    \end{tabular}

    \subsection*{Pérdidas Ópticas}
    \begin{tabular}{|l|c|}
    \hline
    \textbf{Elemento} & \textbf{Pérdida} \\
    \hline
    Fibra SMF @ 1550nm & 0.18-0.20 dB/km \\
    Conector & 0.3-0.5 dB \\
    Patch cord & 0.5-0.6 dB \\
    MUX/DEMUX 40ch & 5 dB \\
    MUX/DEMUX 80ch & 6 dB \\
    Splice & 0.05-0.1 dB \\
    \hline
    \end{tabular}

    \subsection*{Amplificadores Ópticos}
    \begin{tabular}{|l|c|c|}
    \hline
    \textbf{Tipo} & \textbf{Ganancia} & \textbf{NF} \\
    \hline
    EDFA & 20-30 dB & 4-6 dB \\
    Raman & 10-20 dB & 3-5 dB \\
    SOA & 10-25 dB & 6-9 dB \\
    \hline
    \end{tabular}
    \textbf{EDFA:} C-band (1530-1565nm) • $G_{max} \approx 30$ dB
    $P_{sat} \approx +17$ dBm • $B_w \approx 35$ nm

    \subsection*{Tipos Fibra}
    \begin{tabular}{|l|c|c|}
    \hline
    \textbf{Tipo} & \textbf{Ganancia} & \textbf{NF} \\
    \hline
    EDFA & 20-30 dB & 4-6 dB \\
    Raman & 10-20 dB & 3-5 dB \\
    SOA & 10-25 dB & 6-9 dB \\
    \hline
    \end{tabular}
    \textbf{EDFA:} C-band (1530-1565nm) • $G_{max} \approx 30$ dB
    $P_{sat} \approx +17$ dBm • $B_w \approx 35$ nm

    \subsection*{Tipos Fibra}
    \textbf{SMF G.652:} $\alpha$ = 0.18-0.20 dB/km • $D$ = 17 ps/(nm·km) • Núcleo 9μm
    \textbf{NZDSF G.655:} $D$ = 2-6 ps/(nm·km)
    \textbf{DCF:} $D$ = -80 a -100 ps/(nm·km)
    \textbf{MMF:} Núcleo 50/62.5μm • $\alpha$ = 2-3 dB/km • < 2km

    \subsection*{Órbitas Satélite}
    \textbf{LEO:} 500-2000 km • Retardo 3-10 ms
    \textbf{MEO:} 2000-35786 km • GPS @ 20200km
    \textbf{GEO:} 35786 km • Retardo 250-280 ms

    \subsection*{Tasas Datos Estándar}
    \textbf{SDH/SONET:} STM-1/OC-3 (155 Mb/s), STM-4/OC-12 (622 Mb/s)
    STM-16/OC-48 (2.5 Gb/s), STM-64/OC-192 (10 Gb/s)
    \textbf{Ethernet:} 1GbE, 10GbE, 40GbE, 100GbE, 400GbE
    \textbf{OTN:} OTU1 (2.7G), OTU2 (10.7G), OTU3 (43G), OTU4 (112G)

    \section*{CONCEPTOS CLAVE}

\subsection*{Temperatura Ruido}
$T_{sys}$ engloba TODAS fuentes ruido sistema:
$$T_{sys} = T_{antena} + T_{receptor}$$
\textbf{Componentes:}
• $T_{antena}$: ruido captado del entorno
  - Cielo despejado: 10-50 K (baja emisión térmica)
  - Lluvia/nubes: 100-290 K (absorción atm. → emisión)
  - Tierra/obstáculos: 290 K (temperatura física ambiente)
• $T_{receptor} = T_{eq} = T_0(F-1)$: ruido generado por equipo
  - Depende de NF del LNA (primera etapa crítica)
  - Típico: 50-150 K para LNA de calidad
\textbf{Teoría:}
• Teorema Nyquist: resistencias generan ruido térmico
• Menor $T_{sys}$ = mejor sensibilidad receptor
• Sistemas satelitales: $T_{sys}$ bajo crítico (señales débiles)

\subsection*{Link Budget}
\textbf{Objetivo:} Garantizar $C/N$ o $E_b/N_0$ suficiente para BER requerida
\textbf{Metodología:}
1. \textbf{Potencia recibida:}
   $P_R = \text{EIRP} + G_{RX} - A_{total}$ [dB]
   • EIRP = $P_{TX} + G_{TX}$ (potencia radiada efectiva)
   • $A_{total}$: FSL + atm + lluvia + apuntamiento + polarización
2. \textbf{Potencia ruido:}
   $N = kT_{sys}B$ calculado en canal receptor
   • $B$: ancho banda receptor (filtro IF o equivalente)
3. \textbf{Relación portadora/ruido:}
   $C/N = P_R - N$ [dB] → comparar con umbral modulación
4. \textbf{Margen sistema:}
   $M = (C/N)_{real} - (C/N)_{req} \geq 3-6$ dB
   • Compensa variaciones propagación, envejecimiento, errores
\textbf{Concepto:}
• Balance energético completo: TX → Canal → RX
• Todos términos en dB: sumas/restas (NO multiplicaciones)

\subsection*{Multi-hop y Cascada}
$$\frac{1}{(C/N)_{total}} = \sum_i \frac{1}{(C/N)_i}$$ [VALORES LINEALES]
\textbf{Concepto:}
• Ruidos independientes se suman en POTENCIA (lineal)
• Sistema total PEOR que eslabón más débil
• Aplicaciones:
  - Enlace satelital: uplink + downlink
  - Sistemas multi-salto (repetidores)
  - Combinación interferencias múltiples
\textbf{Friis (cascada componentes):}
• Primera etapa domina ruido total sistema
• LNA con bajo NF y alta ganancia minimiza degradación
• Razón: ganancia $G_1$ divide contribución etapas siguientes
\textbf{Estrategia diseño:}
• Optimizar primer eslabón (LNA, primer amp)
• Eslabón débil limita prestaciones globales

\subsection*{EDFA y WDM}
\textbf{Principio:}
• EDFA amplifica TODOS canales WDM simultáneamente (transparente)
• Amplificación en banda (C-band: 1530-1565 nm)
\textbf{Relaciones potencia:}
$$P_{total} = P_{canal} + 10\log_{10}(N_{ch})$$ [dBm]
• Potencia dividida entre canales: $P_{ch} = P_{out} - 10\log(N_{ch})$
• Más canales → menos potencia por canal → peor OSNR
\textbf{Degradación OSNR:}
• Cada EDFA añade ASE (ruido acumulativo)
• OSNR degrada: $\Delta\text{OSNR} = 10\log(N_{amp})$ [dB]
• Límite práctico: OSNR mínimo para modulación (15-25 dB)
\textbf{Diseño enlace:}
$$N_{amp} = \lceil L/D_{amp}\rceil - 1$$ donde $D_{amp} = G_{EDFA}/\alpha$
• $D_{amp}$: distancia entre amplificadores (típ. 80-100 km)
• Redondear hacia ARRIBA, luego restar 1 (sin amp al final)

\section*{PROCEDIMIENTOS}

    \subsection*{Link Budget RF}
    1. Calcular $G_{TX}$ y $G_{RX}$ con fórmula antena
    2. $A_{FSL}$ según distancia/frecuencia (92.45 para GHz)
    3. Sumar pérdidas adicionales ($A_{atm}$, $A_{rain}$, $A_{pol}$, $A_{point}$)
    4. $P_R = P_{TX} + G_{TX} + G_{RX} - A_{total}$
    5. Calcular $N = -174 + 10\log(B) + 10\log(T_{sys}/290)$
    6. $C/N = P_R - N$ • Comparar con umbral modulación
    7. Multi-hop: Sumar $1/(C/N)_i$ en LINEAL, convertir a dB

    \subsection*{Enlace Óptico}
    1. Calcular $N_{amp}$ según $L$ y $D_{amp}$ (redondear arriba -1)
    2. Determinar $N_{ch}$ según $B_{amp}$ y espaciado ITU-T
    3. $P_{ch} = P_{out,EDFA} - 10\log(N_{ch})$
    4. OSNR: $P_{ch} - F_{dB} - 58 - 10\log(N_{amp})$
    5. Comparar OSNR con requisito modulación
    6. Verificar dispersión: $D_{tot} = D \times L$ < límite
    7. Sin EDFA: $L_{max} = (P_{TX} - P_{RX,sens} - L_{conn} - L_{MUX})/\alpha$

    \subsection*{Enlace Satelital Específico}
    \textbf{Uplink:} $(C/N)_{up} = \text{EIRP}_{ES} + G/T_{sat} - A_{up} - k_{dB}$
    \textbf{Downlink:} $(C/N)_{down} = \text{EIRP}_{sat} + G/T_{ES} - A_{down} - k_{dB}$
    $k_{dB} = 10\log(1.38\times10^{-23}) = -228.6$ dBW/K/Hz
    \textbf{Total:} $1/(C/N)_{tot} = 1/(C/N)_{up} + 1/(C/N)_{down}$ [LINEAL]
    \textbf{Distancia GEO:} $d \approx 35786 + 6371 - 6371\cos\varepsilon \approx 36000$-42000 km

    \end{multicols}

    \end{document}
